\section{What the actual PW tail estimate proves for every compact test}
\label{sec:all-test-pw-tail}

\paragraph{MA-AT-001 (ambient trace class and a uniform tail, without a
negative-time divisor formula).}
Fix \(0<\alpha<1/2\) and \(s\ge0\), and use the one physical Hilbert space
\[
 H=L^2\!\left(\mathbb R,w(x)\,dx\right),\qquad
 w(x)=e^{2\alpha|x|}(1+x^2)^{s/2}.
 \tag{MA-AT.1}
\]
For \(c\le0\), let \(P_c=M_{\mathbf1_{[c,\infty)}}\), let \(J_c\) be the
actual closed global PW range with support edge \(c\), and let
\[
 Q_c=P_{J_c},\qquad P_{Z_c}=P_c-Q_c,\qquad
 P_+=P_0.
 \tag{MA-AT.2}
\]
Thus \(J_c\subset P_cH\), and \(J_0\subset J_c\).
The range assertion is the actual one of MA-GP-002 and MA-GW-001,
not a replacement of a source norm by an ordinary Hardy norm.
The same lemma applies to another actual nested PW family whenever
its causal ambient lift at each fixed \(c\) has been proved.

For every \(g\in C_c^\infty(\mathbb R)\), without a sign restriction
on its support,
\[
 C_gP_{Z_c}\in\mathcal S_1(H),\qquad
 \tau_c(g):=\operatorname{Tr}_H(C_gP_{Z_c})
 \quad(c\le0).
 \tag{MA-AT.3}
\]
This defines an actual ambient trace.  For a noncausal test it is
not being identified here with a convolution action on a quotient
that only carries positive translations.

For a nonzero test put
\(a=\min\operatorname{supp}g\) and
\(b=\max\operatorname{supp}g\).
For \(R>b\) there is the exact factorization
\[
 P_{[R,\infty)}C_gP_{Z_c}
   =P_{[R,\infty)}C_gP_{Z_0}P_+P_{Z_c}.
 \tag{MA-AT.4}
\]
Consequently the fixed tail function
\[
 \omega_g(R)=\|P_{[R,\infty)}C_gP_{Z_0}\|_1
       \longrightarrow0\quad(R\longrightarrow\infty)
 \tag{MA-AT.5}
\]
satisfies
\[
 \sup_{c\le0}\|P_{[R,\infty)}C_gP_{Z_c}\|_1
       \le\omega_g(R)\qquad(R>b).
 \tag{MA-AT.6}
\]
All projections and norms in these formulas use the same \(H\).
No equivalence constant depending on \(c\) is used.

Here is the precise consequence for the actual source cutoff.
Write \(Y=e^{-c}\ge1\), let \(0\le\rho_{X,\delta}\le1\) be a smooth
function supported on \([1/X,X]\), equal to one on
\([1/X+\delta,X-\delta]\), and assume
\[
 X\ge 2\max(1,e^{-a})Y,\qquad
 0<\delta\le(4X)^{-1},\qquad \log(X-\delta)>b.
 \tag{MA-AT.7}
\]
Then every trace in
\[
 \begin{split}
 \bigg|&
 \operatorname{Tr}_H(M_{\rho_{X,\delta}(e^\cdot)}C_gQ_c)
       -g(0)\int_c^\infty\rho_{X,\delta}(e^x)\,dx
       +\tau_c(g)\bigg|\\
 &\hspace{20mm}\le
       \omega_g\bigl(\log(X-\delta)\bigr)
 \end{split}
 \tag{MA-AT.8}
\]
is the trace of an actual trace-class operator.
In particular, along every source-thin path with \(Y/X\to0\),
the left expression tends to zero.  If \(g(0)=0\), the displayed
volume counterterm is absent.  The result reduces the all-test
joint-limit problem to the separate question of the limit of
the actual numbers \(\tau_{-\log Y}(g)\).

\begin{proof}
\emph{Translation gives the ambient lift for an arbitrary test.}
Let \(U_t f(x)=f(x-t)\) on the full space \(H\).
Moderateness of \(w\) gives, for every fixed \(t\),
\[
 \|U_t\|\le 2^{s/4}e^{\alpha|t|}(1+t^2)^{s/4}.
 \tag{MA-AT.9}
\]
Indeed \(e^{2\alpha|x+t|}\le
e^{2\alpha|t|}e^{2\alpha|x|}\) and
\(1+(x+t)^2\le2(1+t^2)(1+x^2)\).
Choose \(t_0>-a\) and set \(g_{t_0}(t)=g(t-t_0)\).
Then \(g_{t_0}\in C_c^\infty((0,\infty))\) and
\[
 C_{g_{t_0}}=U_{t_0}C_g,\qquad
 C_gP_{Z_c}=U_{-t_0}C_{g_{t_0}}P_{Z_c}.
 \tag{MA-AT.10}
\]
MA-GW-003 proves the second factor on the right is trace class
on this actual \(H\), for each fixed \(c\).
The first factor is bounded.  This proves (MA-AT.3), without
commuting a translation past \(P_{Z_c}\) and without constructing
a negative-time action on \(H/J_c\).

\emph{Nesting and a physical support calculation.}
If \(v\in H\) and \(j\in J_0\), then \(P_+j=j\) and \(j\in J_c\).
Orthogonality in \(H\) therefore gives
\[
 \langle P_+P_{Z_c}v,j\rangle_H
     =\langle P_{Z_c}v,j\rangle_H=0.
 \]
The vector \(P_+P_{Z_c}v\) is supported on \([0,\infty)\),
so it belongs to \(Z_0=P_+H\ominus J_0\).  Hence
\[
 P_+P_{Z_c}=P_{Z_0}P_+P_{Z_c},\qquad
 \|P_+P_{Z_c}\|\le1.
 \tag{MA-AT.11}
\]
There is no assertion that the two projections commute.
For \(x\ge R>b\), the integral
\(C_gf(x)=\int g(t)f(x-t)\,dt\) samples only \(x-t>0\).
Consequently \(P_{[R,\infty)}C_g
 =P_{[R,\infty)}C_gP_+\), first on compactly supported
vectors and then by boundedness on \(H\).
Combining this with (MA-AT.11) proves (MA-AT.4).
Since \(C_gP_{Z_0}\) is a fixed trace-class operator,
strong convergence \(P_{[R,\infty)}\to0\) implies (MA-AT.5).
The contraction in (MA-AT.11) proves the uniform estimate.

\emph{The actual cutoff trace and its diagonal term.}
For fixed compact position support, the operator
\(M_{\rho(e^\cdot)}C_gP_c\) is trace class.
To see this in the actual weighted space, conjugate by multiplication
by \(\sqrt w\).  Compact support of \(\rho\) and \(g\) confines both
kernel variables to a compact rectangle.  The conjugated kernel is
a smooth compact convolution kernel multiplied on the two sides by
bounded compactly supported functions, together with the bounded
input projection \(P_c\).  The smooth kernel is trace class, for
example by its rapidly summable Fourier coefficients on a containing
rectangle.  Thus the nonsmooth but bounded weight factors and the
sharp support projection cause no trace-class problem.
Its trace is
\[
 \operatorname{Tr}_H(M_{\rho(e^\cdot)}C_gP_c)
       =g(0)\int_c^\infty\rho(e^x)\,dx.
 \tag{MA-AT.12}
\]
For precision, cyclically move the bounded multiplication factors
against that fixed smooth trace-class kernel, so their product on
the diagonal is \(\rho(e^x)\mathbf1_{[c,\infty)}(x)\);
the two weight factors cancel.
The usual smooth-kernel diagonal formula then extends to this
bounded compact multiplier by bounded smooth approximation,
strong multiplier convergence and trace-ideal continuity.
This proves (MA-AT.12) also at the sharp endpoint \(c\).

The identity \(Q_c=P_c-P_{Z_c}\), (MA-AT.3), and the preceding
compact-kernel argument show that the first operator in (MA-AT.8)
is trace class and that the expression before its absolute value is
\[
 \operatorname{Tr}_H
     \bigl((1-M_{\rho_{X,\delta}(e^\cdot)})C_gP_{Z_c}\bigr).
 \tag{MA-AT.13}
\]
The image of \(C_gP_{Z_c}\) is supported on \([c+a,\infty)\).
By (MA-AT.7),
\[
 1/X+\delta\le5/(4X)<e^a/Y=e^{c+a},
 \]
so the entire lower transition is below that image.
The multiplier in (MA-AT.13) can therefore be restricted to
\([R,\infty)\), where \(R=\log(X-\delta)>b\).
It has norm at most one.  Apply (MA-AT.6) and
\(|\operatorname{Tr}T|\le\|T\|_1\) to obtain (MA-AT.8).
\end{proof}

\paragraph{Exact boundary of this result.}
For a causal test, MA-GW-002--003 identify \(\tau_c(g)\) with a
fixed half-plane divisor trace, independent of \(c\), and the
all-weight causal slow-band theorem follows.
For a test with negative support, (MA-AT.3)--(MA-AT.8) establish
actual trace class and a uniform cutoff error, but do not provide
a divisor formula for \(\tau_c(g)\) or its limit as \(c\to-\infty\).
In particular, a bounded full-space translation used in
(MA-AT.10) is not an inverse for the positive semigroup on the
PW quotient.  No cyclic trace step can replace that missing
identification.
The finite-length spectral modules used in MA-GW describe
positive convolution and its trace.  This lemma is an analytic
statement about their actual ambient completions and orthogonal
projections; their class in the trace framework does not supply
the unproved negative-time formula.
